\documentclass[10pt,a4paper]{amsart}
\usepackage[margin=19mm]{geometry}
\usepackage{amsmath,amssymb,mathtools}
\usepackage[scr=rsfs,cal=euler]{mathalfa}
\usepackage{xfrac}
\usepackage[unicode,hidelinks,bookmarks=false]{hyperref}
\newcommand{\ie}{i.e.\ }
\newcommand{\cf}{cf.\ }
\newcommand{\resp}{respectively\ }
\newcommand{\Subsection}{\subsection}
\providecommand{\nobreakdash}{\nobreak\hbox{-}}
\setlength{\emergencystretch}{2em}
\multlinegap0pt
\usepackage{bm}
\usepackage{bbm}
\usepackage{dsfont}
\usepackage{xspace}
\usepackage{cancel}
\usepackage{mathtools}
\usepackage{amsthm}
\makeatletter
\@ifundefined{thm}{\newtheorem{thm}{Thm}}{}
\@ifundefined{prop}{\newtheorem{prop}[thm]{Proposition}}{}
\makeatother
\title[Constructing far-from-equilibrium patterns in a cross-diffus]{Constructing far-from-equilibrium patterns in a cross-diffusion vegetation-autotoxicity model}
\author{Annalisa Iuorio, Cinzia Soresina, Frits Veerman}
\date{Source: arXiv:2607.15692v1 (CC BY 4.0)}
\begin{document}
\maketitle

\noindent\emph{Source and licence.} Constructing far-from-equilibrium patterns in a cross-diffusion vegetation-autotoxicity model. Version arXiv:2607.15692v1 (CC BY 4.0), distributed under the Creative Commons Attribution 4.0 International licence, as recorded in the arXiv licence metadata for this exact version and shown on the abstract page https://arxiv.org/abs/2607.15692v1. Full rights notice: licenses/arxiv-2607.15692-CC-BY-4.0.txt.\par\smallskip

\noindent\textit{Editorial scope.} Counted unit: the Prop whose source label is \texttt{prop:hetconn} in the article (source0.tex lines 450--452, source label \texttt{prop:hetconn}), delivered together with the article's own complete original proof of it (source0.tex lines 453--529, one \texttt{proof} environment of 77 lines). No mathematics is written, repaired or re-derived: statement and proof are the article's own text line for line. Required context retained uncounted: eq:gammas0\_layp (lines 412--419); eq:equilCM (lines 424--426). Every definition, display and label of those lines is retained unchanged. Omitted from this package and not redistributed: 1--411, 420--423, 427--449, 530--1241. Nothing inside a retained excerpt is omitted or altered: every retained body is delivered exactly as the article's lines are written, so no omission receipt of any kind accompanies this package. Citations: the retained excerpts carry no citation command at all, so no bibliography entry of the article is transcribed and no reference is redistributed. Reference adaptation: none was needed and none was made; every reference inside the delivered unit resolves inside it, so no unresolved reference or citation is produced. Printed slips kept literal: no printed word, symbol, hypothesis or number of the counted statement or of its proof was altered. Layout and class: the article's own class is \texttt{scrartcl} with packages including babel, amsmath, amssymb, wasysym, amsfonts, amsthm, empheq, bm, graphicx, epsfig, subfig, overpic, txfonts, fontenc; the wrapper is the lane's pinned \texttt{amsart} editorial wrapper at 10pt on A4, which restates the theorem-like environments the retained text needs and adds the article's own macro definitions that the retained text needs (none), with extra wrapper packages (bm, bbm, dsfont, xspace, cancel, mathtools). The delivered numbering is the unit's own. Assets: no retained span contains a figure environment, a graphics-inclusion command, a PDF-page inclusion, a table or a TikZ picture; no image file is copied into this package, the package declares no asset dependency and no image file is redistributed with it. Licence: version arXiv:2607.15692v1 of this article is distributed under the Creative Commons Attribution 4.0 International licence, as recorded in the arXiv licence metadata for this exact version and shown on the abstract page https://arxiv.org/abs/2607.15692v1. A full rights notice is delivered at licenses/arxiv-2607.15692-CC-BY-4.0.txt. Characters: every retained excerpt is pure ASCII, so the delivered engine, XeLaTeX with its default Latin Modern text font, reports no missing character.
\begin{subequations}\label{eq:gammas0_layp}
    \begin{align}
        % \dot{P} &= 0,\\
        % \dot{S} &= 0,\\
        \dot{T} &= U,\\
        \dot{U} &= T - \frac{P}{1 - \rho\,\phi(T)},
    \end{align}
\end{subequations}

\begin{equation} \label{eq:equilCM}
    T_1^\ast = \frac{1-\sqrt{1-4\rho \bar{P}}}{2 \rho}, \qquad T_2^\ast = \frac{1+\sqrt{1-4\rho \bar{P}}}{2 \rho}, \qquad T_3^\ast = \frac{\bar{P}}{1-\rho}.
\end{equation}

\begin{prop} \label{prop:hetconn}
    For any $1/2 < \rho < 1$ and any $S=\bar{S} \in \mathbb{R}$, there exists a unique $P^\ast(\rho)$ with $1-\rho < P^\ast(\rho) < 1/(4\rho)$ such that system \eqref{eq:gammas0_layp} admits a double heteroclinic connection $\varphi_\mathrm{h}^\pm(\xi) = \left(T_\mathrm{h}(\xi),\pm U_\mathrm{h}(\xi)\right)$ between $Q_1=(T_1^\ast, 0 )$ and $Q_3=(T_3^\ast, 0 )$. In other words, for $P = P^\ast(\rho)$, the intersection $\mathcal{W}^{u}(Q_1) \cap \mathcal{W}^{s}(Q_3)$ (resp.~$\mathcal{W}^{s}(Q_1) \cap \mathcal{W}^{u}(Q_3)$) is transversal.
\end{prop}

\begin{proof}
% Our strategy is the following:
% \begin{itemize}
%     \item Use Hamiltonian structure of the layer problem on $\mathcal{C}_0^l$ and $\mathcal{C}_0^r$ to find orbits $U_{1,3}(T)$ passing through $Q_1$ and $Q_3$ (one corresponding to $\mathcal{W}^s$, one to $\mathcal{W}^u$)
%     \item Prove that the intersection $\mathcal{W}^{s,u}(Q_1) \cap \mathcal{W}^{u,s}(Q_3)$ is transversal, i.e.~for any $\rho$ there exists a unique $\bar{P}(\rho)$ with $1-\rho < \bar{P} < \frac{1}{4\rho}$ such that a double heteroclinic connection between $Q_1$ and $Q_3$ exists.
% \end{itemize}
The vector field of the layer problem \eqref{eq:gammas0_layp} is non-smooth for $T_1^\ast < T < T_3^\ast$, since $1 \in (T_1^\ast, T_3^\ast)$. In particular, at $T=1$, the vector field is continuous but not continuously differentiable. To determine the shape of the stable and unstable manifolds of $Q_1$ and $Q_3$, we study the vector field on both sides of the vertical line $T=1$.\\
The orbits of system \eqref{eq:gammas0_layp} passing through points on $\mathcal{C}_0^l$ are obtained by solving
\begin{equation}
\begin{aligned}
    \frac{\text{d}U}{\text{d}T} &= \frac{1}{U} \left(T - \frac{P}{1 - \rho\,T} \right), \\
    U(T_1^\ast) &= 0,
\end{aligned}
\end{equation}
by separation of variables, which leads to
\begin{equation}
    U_1^\pm (T) = 
    \pm \sqrt{T^2 - (T_1^\ast)^2 + 2 \frac{P}{\rho} \log \frac{1-\rho T}{1-\rho T_1^\ast}},\quad \frac{1-\rho}{\rho} < T < 1.
    %\pm \frac{\sqrt{2 \rho P +\sqrt{1-4 \rho P }-4 \rho P  \log \left(\frac{1}{2} \left(\sqrt{1-4 \rho P }+1\right)\right)+4 \rho P  \log (1-\rho  T)+2 \rho ^2 T^2-1}}{\sqrt{2} \rho }.
\end{equation}
In particular, we have that the stable (resp.~unstable) manifold $\mathcal{W}^{s}$ (resp.~$\mathcal{W}^{u}$) of a generic point $Q_1 \in \mathcal{C}_0^l$ to the left of the vertical line $T=1$ is given by
\begin{equation}
\begin{aligned}
    \mathcal{W}^{s} (Q_1) &= \left\{ (P, S, T, U) \, : \, P=\bar{P}, \, S=\bar{S},\, \frac{1-\rho}{\rho}<T<1, \,  U=U_1^-(T)  \right\}, \\
    \mathcal{W}^{u} (Q_1) &= \left\{ (P, S, T, U) \, : \, P=\bar{P}, \, S=\bar{S},\, \frac{1-\rho}{\rho}<T<1, \,  U=U_1^+(T)  \right\}.
\end{aligned}
\end{equation}
Analogously, the orbits of system \eqref{eq:gammas0_layp} passing through points on $\mathcal{C}_0^r$ are obtained by solving
\begin{equation} \label{eq:UTsys_3}
\begin{aligned}
    \frac{\text{d}U}{\text{d}T} &= \frac{1}{U} \left(T - \frac{P}{1 - \rho} \right), \\
    U(T_3^\ast) &= 0,
\end{aligned}
\end{equation}
yielding
\begin{equation}
    U_3^\pm (T) = \pm \left(\frac{P}{1-\rho } - T \right),\quad 1 < T < \frac{1}{4 \rho  (1-\rho )}.
\end{equation}
In this case, because of the linearity of system \eqref{eq:UTsys_3} when applying the method of separation of variables, stable and unstable manifolds of $Q_3 \in \mathcal{C}_0^r$ coincide on the interval $T \in \left(1, 1/(4 \rho  (1-\rho ))\right)$ with the corresponding stable and unstable eigenspaces $E^{s,u}(Q_3)$, namely
\begin{equation}
\begin{aligned}
    \mathcal{W}^{s} (Q_3) &= \left\{ (P, S, T, U) \, : \, P=\bar{P}, \, S=\bar{S},\, 1<T<\frac{1}{4 \rho  (1-\rho )}, \,  U=U_3^-(T)  \right\} = E^{s} (Q_3), \\
    \mathcal{W}^{u} (Q_3) &= \left\{ (P, S, T, U) \, : \, P=\bar{P}, \, S=\bar{S},\, 1<T<\frac{1}{4 \rho  (1-\rho )}, \,  U=U_3^+(T)  \right\} = E^{u} (Q_3).
\end{aligned}
\end{equation}
Our goal is now to prove that the intersection $\mathcal{W}^{s}(Q_1) \cap \mathcal{W}^{u}(Q_3)$ (resp.~$\mathcal{W}^{u}(Q_1) \cap \mathcal{W}^{s}(Q_3)$) is transversal. Thanks to the symmetry of the functions $U_{1,3}^\pm$ involved, it suffices to prove this result for one intersection, and the other one will follow directly. Therefore, we focus on $\mathcal{W}^{u}(Q_1) \cap \mathcal{W}^{s}(Q_3)$.\\
We check the intersection of these manifolds at $T=1$. 
%\footnote{\textcolor{red}{Because of piecewise $\phi$, we should check that the resulting orbit is sufficiently regular.}}. 
That is, our aim is to show that there exists a unique solution $\bar{P}=P^\ast(\rho)$ with $1-\rho < P^\ast(\rho) < 1/(4\rho)$ such that $U_1^+(1)=U_3^-(1)$, i.e.~$U_1^+(1)-U_3^-(1)=0$. 
%This derives directly from the properties of the functions involved; in particular
In order to show the existence and uniqueness of $\bar{P}=P^\ast(\rho)$, we note the following:
\begin{itemize}
    \item $U_1^+(1)$ is a strictly monotonically decreasing function of $\bar{P}$, since
    \begin{displaymath}
        \frac{\text{d}}{\text{d} P}U_1^+(1) = \frac{1}{\rho U_1^+(1)} \log \frac{1-\rho}{1-\rho T_1^\ast} < 0 
    \end{displaymath}
    because $T_1^\ast<1$. Moreover, $U_1^+(1)$ assumes
    a positive value at $\bar{P}=1-\rho$ (namely, the lower extremum of the range of admissible $\bar{P}$-values) and admits a zero at a value of $\bar{P} < 1/(4\rho)$.
    %\textcolor{red}{FV: I would add a calculation here showing that there is a $\bar{P}$-value for which $U_1^+(1) = 0$, or at least $U_1^+(1) \downarrow 0$ for $\bar{P} \uparrow \bar{P}_*$}
    %\footnote{\textcolor{red}{Show details here/in an appendix? Or not at all?}}.
    In fact, when substituting $T_1^\ast$ as in Eq.~\eqref{eq:equilCM}, we have that
    \[
    U_1^+(1) = \frac{\sqrt{2 \rho (\bar{P} + \rho) +\sqrt{1-4 \rho \bar{P} }+4 \rho \bar{P}  \log \left(\dfrac{2(1- \rho) }{\sqrt{1-4 \rho \bar{P} }+1}\right)-1}}{\sqrt{2} \rho },
    \]
    which satisfies $U_1^+(1) = 0 $ when
    \[
    2 \rho (\bar{P} + \rho) +\sqrt{1-4 \rho \bar{P} }+4 \rho \bar{P}  \log \left(\frac{2(1- \rho) }{1+\sqrt{1-4 \rho \bar{P} }}\right)-1=0,
    \]
    or, equivalently, when
    \[
    u_1^+(\bar{P}):=e^{-(2 \rho  (\bar{P}+\rho )+\sqrt{1-4 \rho \bar{P} }-1)/4 \rho \bar{P} }-\frac{2 (1-\rho )}{1+\sqrt{1-4 \rho \bar{P} }} = 0.
    \]
    The function $u_1^+(\bar{P})$ is negative at $\bar{P}=1-\rho$ and positive at $\bar{P}=1/(4 \rho)$ for any $1/2 < \rho < 1$; therefore, by the intermediate value theorem, it must admit a zero within this $\bar{P}$-range. 
    \item $U_3^-(1)$ is a linear, strictly monotonically increasing function of $\bar{P}$, equal to $0$ at $\bar{P}=1-\rho$.
\end{itemize}
Therefore, for any $\rho \in \left(1/2, 1 \right)$ there exists a unique $\bar{P}=P^\ast(\rho)$ with $1-\rho < P^\ast(\rho) < 1/(4\rho)$ such that $(U_1^+(1)-U_3^-(1))|_{\bar{P}(\rho)} = 0$. The transversality of the intersection $\mathcal{W}^{s}(Q_1) \cap \mathcal{W}^{u}(Q_3)$ follows from the non-degeneracy of $\bar{P}(\rho)$, i.e. from the transversality of the intersection of $U_1^+(1)$ and $U_3^-(1)$ as graphs over $P$.
\end{proof}

\end{document}
